\documentclass[conference]{IEEEtran}
\IEEEoverridecommandlockouts
\usepackage{cite}
\usepackage{amsmath,amssymb,amsfonts}
\usepackage{algorithmic}
\usepackage{graphicx}
\usepackage{textcomp}
\usepackage{xcolor}
\usepackage[hidelinks]{hyperref}
\def\BibTeX{{\rm B\kern-.05em{\sc i\kern-.025em b}\kern-.08em
    T\kern-.1667em\lower.7ex\hbox{E}\kern-.125emX}}

\begin{document}

\title{Unsupervised Neural Network for Multi-Genre Music Generation Using LSTM Autoencoder, VAE, Transformer, and Reinforcement Learning}

\author{
\IEEEauthorblockN{
1\textsuperscript{st} Md. Farhan Mahtab Sohan \\
\textit{ID: 22301272}
}
\IEEEauthorblockA{\textit{Dept. of CSE} \\
\textit{BRAC University}\\
Dhaka, Bangladesh \\
farhan.mahtab.sohan@g.bracu.ac.bd}
\and
\IEEEauthorblockN{
2\textsuperscript{nd} Fariha Rahman \\
\textit{ID: 22301307}
}
\IEEEauthorblockA{\textit{Dept. of CSE} \\
\textit{BRAC University}\\
Dhaka, Bangladesh \\
fariha.rahman4@g.bracu.ac.bd}
}

\maketitle

% ============================================================
\begin{abstract}
Automated generation of music is a difficult issue as there are complex temporal issues in music.
structure and genre-specific stylistic patterns intrinsic to musical signals.
The complexity of acquiring supervised methods of music generation is expensive.
labeled musical data. In this paper, our hypothesis is a progressive unsupervised learning.
generative model to produce and create multi-genre music without the use of explicit.
genre labels in training. Four generative models of our system are developed.
enhancing capacity: an LSTM Autoencoder to reconstruct single-genre music, a
A Transformer called Variational Autoencoder (VAE) of diverse multi-genre generation.
long coherent sequence synthesis decoder, and a Reinforcement Learning of.
Preference-aligned generation Human Feedback (RLHF) stage. We train and
test our models with two publicly available MIDI datasets: the MAESTRO dataset,
including the high-quality Classical piano performances, and the Lakh MIDI Dataset,
spanning a wide spectrum of genres such as Jazz, Rock, Pop and Electronic.
Measurement is done based on quantitative measures such as pitch histogram.
distance, diversity score of rhythm, ratio of repetitions, and perplexity. The
Transformer has an ultimate validation perplexity of 3.8447 and optimal.
validation loss of 1.3467, which significantly does better than Markov Chain and random.
baseline models. The RLHF-tuned model does not decrease the Transformer base perplexity.
of 3.8447 and increasing the means of human listening scores of 3.7056 to 3.9941.
on 5-point Likert scale among 18 respondents, affirming that human.
preference alignment generates uniformly superior quality music without.
worsening the token-level prediction quality of the model. These results demonstrate
that unconditioned generative models have the ability to generate musically coherent and
genre-consistent non-manually annotated compositions.
\end{abstract}

\begin{IEEEkeywords}
music generation, MIDI, LSTM autoencoder, variational autoencoder, transformer,
reinforcement learning from human feedback, unsupervised learning, deep generative models
\end{IEEEkeywords}


% ============================================================
\section{Introduction}

Music is a highly organized time cue that is used to encode several concurrent.
patterns, such as melodic progressions, harmonic relationships, rhythmic variations,
and distribution of style over genres. The automatic creation of music that
these patterns have long been a challenge to faithfully capture in the
crossover between signal processing and machine learning.~\cite{briot2017deep}.


Older methods of music generation have been based on rule-based systems and
statistical models like Markov chains which are short-range dependent but
has a problem of musical continuity.~\cite{papadopoulos1999ai}. More recently, deep supervised learning techniques have demonstrated good results in music modeling, yet
they need vast amounts of labelled data which are costly and time consuming.
to obtain. This is an incentive to use unsupervised learning generative models.
extremely high-level musical expressions of raw symbolic MIDI data without explicit.
genre annotations.

Formally, we represent a music sequence as:
\begin{equation}
X = \{x_1, x_2, \ldots, x_T\}
\end{equation}
where each $x_t$ is a symbolic event such as a note-on event, note-off event,
velocity value, or duration. The central objective is to learn an unsupervised
generative distribution $p_\theta(X)$ such that samples drawn from it:
\begin{equation}
\hat{X} \sim p_\theta(X)
\end{equation}
produce realistic, genre-consistent musical compositions.

This paper makes the following contributions:
\begin{itemize}
    \item We create and establish a four-stage progressive model of unsupervised.
    music generation, with a simple LSTM Autoencoder up to an RLHF-trained.
    Transformer.
    \item  We compare each step to two baselines (Random Note Generator and).
    Markov Chain) quantitatively: pitch histogram distance, rhythm
    diversity, the ratio of repetitions, and perplexity.
.
    \item 
    We show that the Transformer, which is trained on 2,794,017 token windows.
    on a vocabulary of 195 tokens, has a best validation perplexity of
    With a smooth convergence rate of 3.8447 and 10 epochs, and RLHF fine-tuning.
    improves mean human listening scores from 3.7056 to 3.9941 (a $7.79\%$
    Improvement relative) of 18 subjects on a 5-point Likert,
    and the RLHF-tuned model has the same base perplexity as that of 3.8447.
\end{itemize}

The remainder of this paper is organized as follows. Section~\ref{sec:related}
reviews related work on neural music generation. Section~\ref{sec:method} describes
our methodology and model architectures. Section~\ref{sec:experiments} presents
experimental results. Section~\ref{sec:discussion} discusses findings and
limitations, and Section~\ref{sec:conclusion} concludes the paper.


% ============================================================
\section{Related Work}
\label{sec:related}

\subsection{RNN and LSTM-Based Music Generation}

Recurrent neural networks especially Long Short-Term Memory (LSTM) networks.
were some of the earliest deep architectures to symbolic music generation.
Eck and Schmidhuber~\cite{eck2002finding} proved the ability of LSTMs to learn.
fundamental blues chords and patterns. Later work e.g.
DeepBach~\cite{hadjeres2017deepbach} and BachBot were built on RNN-based models to
harmonize chorales J.S. Bach style. Regardless of their efficacy in
vanilla RNNs and LSTMs are short sequences that are not very good at preserving coherent structure.
over extended musical phrases because of the vanishing gradient problem.

\subsection{Variational Autoencoders for Music}

VAEs~\cite{kingma2013auto} have been used with music.
to learn continuous, structured latent spaces that can interpolate and control.
cross-fertilization between styles of music. Roberts et MusicVAE model.
A hierarchical latent space was proposed in al. ~\cite{roberts2018hierarchical} and it is stated that.
catches local bar-level structure and global phrase-level structure.
The VAEs enable different sampling using the previous $p(z) = \mathcal{N}(0, I)$,
enabling generation of multi-genres without conditioning of genres.

\subsection{Transformer-Based Music Generation}

The self-attention mechanisms were introduced through the Transformer.
architecture(~\cite{vaswani2017attention} brought new opportunities in capturing.
long-range dependencies of musical sequences. The Music Transformer of Huang et.
Relative position-based attention was introduced in al. ~\cite{huang2018music} in order to model.
the rhythmic form of music. MuseNet by OpenAI and Magenta by Google.
project further showed that massive autoregressive Transformer models have the ability to.
create coherent multi-instrumental compositions of a wide range of genres.

\subsection{Reinforcement Learning for Creative Generation}

Reinforcement Learning with Human Feedback (RLHF) has been proven to be successful in
matching the language model outputs with human preferences~\cite{christiano2017deep}.
The use of this paradigm in music generation is still young. Policy gradient
methods provide a principled method to optimize non-differentiable goals like.
Score human listening scores, which are not directly back-propagable with a.
generative model~\cite{yi2018automatic}. This is where we carry on this research.
through the use of RLHF to generate music in various genres using MIDI.

% ============================================================
\section{Methodology}
\label{sec:method}

\subsection{Data Preprocessing}

We operate MIDI files of two publicly available datasets: the
The MAESTRO dataset~\cite{hawthorne2018enabling} that offers quality data.
Classical music recordings on piano, and the Lakh MIDI.
Dataset~\cite{raffel2016learning} that encompasses a wide multi-genre set.
and Jazz, Rock, Pop, and Electronic styles. The two resources are combined to create.
cover complementary: MAESTRO provides a structurally rich.
single-instrument domain and fine timing, and Lakh MIDI adds genre.
variability in terms of instruments and styles.

MIDI files are all translated to a token-based or piano-roll format.
Timing is normalized to 16 steps per bar and sequences are.
divided into fixed-length chunks of $T$ tokens. The merged data is divided.
into training and test sets in 80/20 proportion.

\subsection{Task 1: LSTM Autoencoder for Single-Genre Generation}

The former is a typical sequence-to-sequence LSTM Autoencoder that is trained on.
the MAESTRO dataset, a clean and structurally consistent single-genre.
(Classical piano) area highly appropriate in teaching an early generative.
representation. The encoder is a map from an input sequence $X$ to a fixed-dimensional latent vector:
\begin{equation}
z = f_\phi(X)
\end{equation}
and the decoder reconstructs the sequence from this latent code:
\begin{equation}
\hat{X} = g_\theta(z)
\end{equation}

Training minimizes the mean squared reconstruction loss:
\begin{equation}
\mathcal{L}_{\text{AE}} = \sum_{t=1}^{T} \|x_t - \hat{x}_t\|^2
\end{equation}

New music is generated by sampling or perturbing the latent vector $z$ and
decoding it through $g_\theta$.

\subsection{Task 2: Variational Autoencoder for Multi-Genre Generation}

The second model is an expansion of the autoencoder to a Variational Autoencoder (VAE).
to facilitate the use of latent space sampling with probability to sample a variety of genres. Training
is extended to cover the entire corpus of MAESTRO (Classical) and Lakh.
MIDI (Jazz, Rock, Pop, Electronic), enabling the model to record a wider.
shared latent space of musical styles. The encoder
parameters of a Gaussian posterior:
\begin{equation}
q_\phi(z|X) = \mathcal{N}(\mu(X),\, \sigma(X))
\end{equation}

A latent sample is drawn using the reparameterization trick:
\begin{equation}
z = \mu + \sigma \odot \epsilon, \quad \epsilon \sim \mathcal{N}(0, I)
\end{equation}

The model is trained with the Evidence Lower Bound (ELBO) objective:
\begin{equation}
\mathcal{L}_{\text{VAE}} = \mathcal{L}_{\text{recon}} + \beta \, D_{\text{KL}}\bigl(q_\phi(z|X) \| p(z)\bigr)
\end{equation}

where the KL divergence is:
\begin{equation}
D_{\text{KL}}(q \| p) = \int q(z) \log \frac{q(z)}{p(z)}\, dz
\end{equation}


The trade-off between reconstruction is regulated by the hyperparameter $\beta$.
fidelity and latent space regularization~\cite{higgins2016beta}. Multi-genre
music is created via drawing $z \sim \mathcal{N}(0, I)$ and decoding.

To further examine the latent space of the VAE we performed a \textbf{latent
interpolation} experiment. Considering two coded vectors $z_1$  and $z_2$ 
latent points between (e.g., Classical and Jazz) are different.
calculated by spherical linear interpolation:
\begin{equation}
z(\alpha) = (1-\alpha)\,z_1 + \alpha\,z_2, \quad \alpha \in [0,1]
\end{equation}
Decoding $z(\alpha)$ at evenly spaced values of $\alpha$ produces compositions
that transition smoothly between source genres. The resulting interpolated MIDI
files are included in the final submission.

\subsection{Task 3: Transformer Decoder for Long-Sequence Generation}

The third model is an autoregressive Transformer decoder that models the
joint probability of a sequence as a product of conditional distributions:
\begin{equation}
p(X) = \prod_{t=1}^{T} p(x_t \mid x_{<t})
\end{equation}

Each token embedding is augmented with a genre embedding:
\begin{equation}
h_t = \text{Emb}(x_t) + \text{Emb}(\text{genre})
\end{equation}

The model is trained by minimizing the negative log-likelihood:
\begin{equation}
\mathcal{L}_{\text{TR}} = -\sum_{t=1}^{T} \log p_\theta(x_t \mid x_{<t})
\end{equation}

Model quality is evaluated using perplexity:
\begin{equation}
\text{Perplexity} = \exp\!\left(\frac{1}{T} \mathcal{L}_{\text{TR}}\right)
\end{equation}

Training runs for 10 epochs on a corpus of 114,714 processed MIDI files
tokenized into a vocabulary of 195 tokens. The training split comprises 91,888
files and 2,794,017 token windows across 559 shards; validation uses 11,532
files and 349,043 windows; the test set contains 11,294 files and 343,553
windows. The cross-entropy loss converges from 1.63 to 1.37 (train) and 1.49
to 1.35 (validation). The corresponding perplexity reduces from 5.09 to 3.93
(train) and from 4.42 to 3.8447 (validation), with a best validation loss of
1.3467. Long compositions are generated by iteratively sampling $x_t \sim
p_\theta(x_t \mid x_{<t})$ token by token.

\subsection{Task 4: Reinforcement Learning from Human Feedback}

The fourth stage fine-tunes the Transformer generator using reinforcement
learning to maximize human listening preference. Music samples are generated
from the current policy:
\begin{equation}
X_{\text{gen}} \sim p_\theta(X)
\end{equation}

Human evaluators assign preference scores:
\begin{equation}
r = \text{HumanScore}(X_{\text{gen}}), \quad r \in [1, 5]
\end{equation}

The expected reward objective is:
\begin{equation}
J(\theta) = \mathbb{E}\bigl[r(X_{\text{gen}})\bigr]
\end{equation}

Parameters are updated via the REINFORCE policy gradient:
\begin{equation}
\nabla_\theta J(\theta) = \mathbb{E}\bigl[r \cdot \nabla_\theta \log p_\theta(X)\bigr]
\end{equation}

A human listening survey was conducted with 18 participants (before RLHF) and
17 participants (after RLHF), who rated 10 generated music samples each on a
5-point Likert scale for overall quality, genre coherence, and creativity. The
reward signal used to drive the policy gradient update was the average human
listening score per sample. A total of 180 ratings were collected before RLHF
and 170 ratings after RLHF, covering 10 scored samples in each condition.
RLHF training ran for 20 RL steps using the REINFORCE policy gradient, with
the average normalized reward and average sequence log-probability monitored
throughout to detect policy collapse. The RLHF-tuned model retains the
Transformer base perplexity of 3.8447, as the policy gradient updates optimize
for human preference rather than token-level likelihood, leaving the base
perplexity unchanged.


% ============================================================
\section{Experiments and Results}
\label{sec:experiments}

\subsection{Experimental Setup}

All models were implemented in Python using PyTorch. Training was conducted on
an NVIDIA GeForce RTX 4080 GPU via Google Colab.

For the Transformer (Task 3), tokenization processed 114,714 MIDI files
(468 skipped) into a vocabulary of 195 tokens. The resulting splits are:
91,888 training files (2,794,017 token windows, 559 shards), 11,532 validation
files (349,043 windows, 70 shards), and 11,294 test files (343,553 windows,
69 shards). Table~\ref{tab:datasets} summarizes the datasets.

\begin{table}[htbp]
\caption{MIDI Datasets Used for Training}
\begin{center}
\begin{tabular}{|l|l|c|}
\hline
\textbf{Dataset} & \textbf{Genre Coverage} & \textbf{Files Used} \\
\hline
MAESTRO   & Classical Piano              & 1276 \\
Lakh MIDI & Jazz, Rock, Pop, Electronic  & 176581 \\
\hline
\textbf{Total} & \textbf{5 Genres} & \textbf{177857} \\
\hline
\end{tabular}
\label{tab:datasets}
\end{center}
\end{table}

\subsection{Baseline Models}

We compare our models against two baselines:
\begin{itemize}
    \item \textbf{Random Note Generator}: Uniformly samples MIDI note events
    at random with no learned structure.
    \item \textbf{Markov Chain Model}: A first-order Markov chain trained on
    note transition probabilities from the training data.
\end{itemize}

\subsection{Evaluation Metrics}

We evaluate all models on four metrics. Pitch Histogram Distance measures
the $L_1$ distance between the pitch distributions of generated and real music:
\begin{equation}
H(p, q) = \sum_{i=1}^{12} |p_i - q_i|
\end{equation}

Rhythm Diversity Score captures the variety of note durations used:
\begin{equation}
D_{\text{rhythm}} = \frac{\text{\# unique durations}}{\text{\# total notes}}
\end{equation}

Repetition Ratio measures structural redundancy:
\begin{equation}
R = \frac{\text{\# repeated patterns}}{\text{\# total patterns}}
\end{equation}

Perplexity evaluates the Transformer's autoregressive prediction quality:
\begin{equation}
\text{Perplexity} = \exp\!\left(\frac{1}{T} \mathcal{L}_{\text{TR}}\right)
\end{equation}

For the RLHF stage (Task 4), model quality is measured by the mean human
listening score $\bar{r}$ aggregated across all participants and samples:
\begin{equation}
\bar{r} = \frac{1}{N \cdot S} \sum_{i=1}^{N} \sum_{j=1}^{S} r_{ij}
\end{equation}
where $N$ is the number of participants and $S$ is the number of scored samples.

\subsection{Quantitative Results}

Table~\ref{tab:results} presents the full comparison of all models. The LSTM
Autoencoder (Task 1) achieves the best pitch fidelity (Pitch Histogram Distance
= 0.1292) but entirely collapses to repetitive outputs (Repetition Ratio =
1.0000, Rhythm Diversity = 0.0114). The VAE (Task 2) substantially improves
Rhythm Diversity to 0.4167 and reduces the Repetition Ratio to 0.1773, at the
cost of a higher Pitch Histogram Distance (0.5868). The Transformer (Task 3)
achieves a best validation perplexity of 3.8447 (best validation loss =
1.3467), a Repetition Ratio of 0.2298, and a Rhythm Diversity of 0.1532, with
a Pitch Histogram Distance of 0.7334. Human listening evaluation was conducted
only in Task 4. The Transformer-generated samples received a mean pre-RLHF
human listening score of 3.7056 ($\sigma = 0.9815$, 18 participants, 180 total
ratings). After RLHF fine-tuning (Task 4), the mean score improved to 3.9941
($\sigma = 0.9671$, 17 participants, 170 total ratings), representing an
absolute improvement of $+0.2886$ points ($7.79\%$ relative improvement). The
RLHF-tuned model retains the Transformer base perplexity of 3.8447, as RLHF
optimizes for human preference rather than token-level likelihood.

% ---------------------------------------------------------------
% TABLE 2 -- CORRECTED from actual output files
% Key corrections vs. original draft:
%   Task 1 Loss: 0.82 -> 0.0339 (MSE, converged train loss from ae_history.json)
%   Task 2 Loss: 0.65 -> 0.0217 (best val ELBO from vae_history.json best_val_total_loss)
%   Task 3 Train PPL (first->last): 5.12->3.94 corrected to 5.09->3.93 (from transformer_history.json)
%   Task 3 Val PPL (first->last):  4.41->3.8447 corrected to 4.42->3.8447
%   Random Baseline Rep. Ratio: 0.0000 (confirmed from metrics_summary.csv)
%   Markov Rep. Ratio: 0.0032 (confirmed from metrics_summary.csv)
% ---------------------------------------------------------------

\begin{table*}[htbp]
\caption{Comprehensive Performance Comparison of Baseline and Proposed Models (Task 1--Task 4)}
\begin{center}
\begin{tabular}{|l|c|c|c|c|c|c|c|c|}
\hline
\textbf{Model} & \textbf{Loss} & \textbf{PPL} & \textbf{Rhy.\ Div.} & \textbf{Rep.\ Ratio} & \textbf{Pitch Dist.} & \textbf{Rhy.\ Level} & \textbf{Human Score} & \textbf{Genre Control} \\
\hline
Random Generator    & --     & --     & 0.0069 & 0.0000 & 0.2090 & Low       & --     & None         \\
Markov Chain        & --     & --     & 0.0968 & 0.0032 & 0.2186 & Medium    & --     & Weak         \\
Task 1: LSTM AE     & 0.0339 & --     & 0.0114 & 1.0000 & 0.1292 & Medium    & --     & Single Genre \\
Task 2: VAE         & 0.0217 & --     & 0.4167 & 0.1773 & 0.5868 & High      & --     & Moderate     \\
Task 3: Transformer & --     & 3.8447 & 0.1532 & 0.2298 & 0.7334 & Very High & 3.7056 & Strong       \\
Task 4: RLHF-Tuned  & --     & 3.8447 & --     & --     & --     & Very High & 3.9941 & Strongest    \\
\hline
\end{tabular}
\label{tab:results}
\end{center}
\noindent\small Note: \textit{Rhy.\ Level} provides a qualitative summary of rhythm diversity.
Task~1 Loss is the converged MSE reconstruction loss (from \texttt{ae\_history.json}).
Task~2 Loss is the best validation ELBO loss (from \texttt{vae\_history.json}), recorded before the KL instability at epoch~18.
Human Score is reported only for Task~3 (pre-RLHF mean = 3.7056, 18 participants, 180 ratings) and Task~4
(post-RLHF mean = 3.9941, 17 participants, 170 ratings); all other models were not subject to human listening evaluation.
PPL for Task~4 is 3.8447, inherited from the Transformer base model, as RLHF policy updates optimize for human
preference rather than token-level likelihood and leave the base perplexity unchanged.
\end{table*}

\subsection{Task 1 Training Curve}

Fig.~\ref{fig:ae_loss} shows the MSE reconstruction loss for the LSTM
Autoencoder over 30 epochs. The training loss drops sharply from approximately
0.0365 at epoch 0 to 0.0339 by epoch 1, then plateaus and remains stable. The
validation loss remains at a relatively low level of 0.0334 across the course of training,
suggesting good generalisation of MAESTRO Classical piano data without.
overfitting.

\begin{figure}[htbp]
\centerline{\includegraphics[width=\columnwidth]{architecture_diagram/ae_loss_curve.png}}
\caption{Task 1 -- LSTM Autoencoder MSE reconstruction loss over 30 epochs.
Train loss (blue) converges rapidly by epoch 1 to 0.0339; validation loss (orange)
remains stable at approximately 0.0334.}
\label{fig:ae_loss}
\end{figure}

In the case of Task 1, the trained LSTM Autoencoder was employed to produce \textbf{5 MIDI
samples} of the Classical genre sampling and perturbing the learned latent.
space. Pitch histogram distance, rhythm were used to analyze these samples.
they contain diversity and ratio of repetition and the generated MIDI files
the last submission as part of the deliverables.

\subsection{Task 2 Training Curves}

The training curves of VAE in 20 epochs across are shown in Fig.~\ref{fig:vae_loss} three subplots: total ELBO loss, reconstruction loss, and KL divergence. The
reconstruction loss increases continuously between $\sim$0.020 and $\sim$0.058 as the
model learns more rich multi-genre representations. There is an acute instability
around epoch 18, the KL divergence peaks at around epoch 18, the KL divergence increases to about $5.9 \times 10^9$
recovers in epoch 20 before recovering. The optimal ELBO of 0.0217 was
this checkpoint is used by recording at epoch 7 before the KL explosion
for generation. This is a KL collapse behaviour under aggressive  scaling by $\beta$ and is further discussed in Section~\ref{sec:discussion}.

\begin{figure}[htbp]
\centerline{\includegraphics[width=\columnwidth]{architecture_diagram/vae_loss_curve.png}}
\caption{Task 2 -- VAE training curves over 20 epochs. Top: total ELBO loss.
Middle: reconstruction loss. Bottom: KL divergence. A transient KL spike
near epoch 18 resolves by epoch 20. Best validation ELBO = 0.0217 (epoch 7).}
\label{fig:vae_loss}
\end{figure}

In Task 2, VAE was created on Classical, Jazz, \textbf{8 MIDI samples}.
Rock, and Pop by sampling $z \sim \mathcal{N}(0, I)$. The latent interpolation
experiment also resulted in \textbf{5 genre-blended compositions}
making a gradual shift between Classical and Jazz; versatility and fluidity
were evaluated through pitch histograms and rhythm diversity measures. All generated
and interpolated MIDI files are part of the final submission.

\subsection{Task 3 Training Curves and Perplexity}

Fig.~\ref{fig:t3_train} demonstrates the Transformer training curves in 10 epochs.
The cross-entropy loss (top) also shows a smooth variation between 1.63 and 1.37.
training set and between 1.49 and 1.35 corresponding to the validation set and the two curves respectively.
meeting closely by epoch 4. The confusion (bottom) decreases respectively.
from 5.09 to 3.93 (train) and 4.42 to 3.8447 (validation), with the best
validation loss recorded 1.3467, which attests to stable autoregressive learning.
without overfitting.

\begin{figure}[htbp]
\centerline{\includegraphics[width=\columnwidth]{architecture_diagram/transformer_training_curves.png}}
\caption{Task 3 -- Transformer training curves over 10 epochs. Top:
cross-entropy loss. Bottom: perplexity. Both train (blue) and validation
(orange) curves converge smoothly, reaching a best validation perplexity
of 3.8447 (loss = 1.3467).}
\label{fig:t3_train}
\end{figure}

The last compared perplexity of the Transformer is reported in Fig.~\ref{fig:t3_ppl}
on the test set. The model has a perplexity of 3.8447 and it is equivalent to.
a minimal validation loss of 1.3467 which validates high next-token prediction.
quality to multi-genre symbolic music generation.

\begin{figure}[htbp]
\centerline{\includegraphics[width=\columnwidth]{architecture_diagram/task3_perplexity.png}}
\caption{Task 3 -- Final Transformer test-set perplexity = 3.8447
(best validation loss = 1.3467).}
\label{fig:t3_ppl}
\end{figure}

In Task 3 the Transformer model has been applied to create the \textbf{10 long-form MIDI
compositions} using autoregressive sampling of tokens by sample in more than one.
genres. Perplexity was used to evaluate these compositions and the results.
establish the ability of the model to generate coherent and long sequence music in general.
Classical, Jazz, Rock, Pop, and Electronic. The MIDI files that are generated are
submitted as part of the final work.

\subsection{Cross-Task Metric Comparison}

The entire metric comparison is shown in Fig.~\ref{fig:all_metrics}
lines and Tasks 1-3 on Pitch histogram distance, Rhythm Diversity,
Repetition Ratio.


\begin{figure}[htbp]
\centerline{\includegraphics[width=\columnwidth]{architecture_diagram/task_comparison_metrics.png}}
\caption{Baseline and task metric comparison for Pitch Histogram Distance
(blue), Rhythm Diversity (orange), and Repetition Ratio (green) across the
random baseline, Markov baseline, Task 1 LSTM Autoencoder, Task 2 VAE,
and Task 3 Transformer.}
\label{fig:all_metrics}
\end{figure}

The lowest Pitch is attained by the LSTM Autoencoder (Task 1).
Histogram Distance (0.1292) of all models except collapses to a
Ratio of Repetition 1.0000 and almost zero Rhythm Diversity
(0.0114). The VAE (Task 2) does not follow this fall, it succeeds
with a Rhythm Diversity of 0.4167 and Repetition Ratio of 0.1773
by stochastic latent sampling. The Transformer
trend is continued in (Task 3) its Repetition Ratio stabilizes
with the 0.2298 and Rhythm Diversity of 0.1532, show more
structured generation-consistent (instead of maximal)
diversity. It has the highest Pitch Histogram Distance (0.7334)
in all models, congruent with the production over a wider
vocabulary of 195 tokens (five-genre). Both baselines score substantially below all learned
models in Rhythm Diversity (Random: 0.0069; Markov: 0.0968), confirming the
advantage of learned generative models.

\subsection{Task 4: RLHF Human Listening Evaluation}

\subsubsection{RLHF Training Dynamics}

Fig.~\ref{fig:rlhf_curve} shows the RLHF training curves over 20 RL update
steps. The RLHF loss (top panel) exhibits high volatility, oscillating between
approximately $-40$ and $+95$ across all steps without a strong downward trend,
which is characteristic of a first-round policy gradient update where the reward
signal is sparse and the policy has not yet stabilized. The average normalized
reward (bottom panel, orange) remains near zero throughout training, while the
average sequence log-probability (blue) stays stable at approximately $-475$,
indicating that the policy did not collapse and remained close to the pretrained
Transformer prior. This log-probability stability is a healthy sign of
constrained policy updates during early-stage RLHF, and is consistent with the
base perplexity of 3.8447 being preserved after fine-tuning.

\begin{figure}[htbp]
\centerline{\includegraphics[width=\columnwidth]{architecture_diagram/task4_rlhf_training_curve.png}}
\caption{Task 4 -- RLHF training dynamics over 20 RL steps. Top: RLHF loss
(high volatility is expected in a first-round policy gradient update). Bottom:
average sequence log-probability (blue, stable near $-475$) and average
normalized reward (orange, near zero), confirming the policy remained close
to the Transformer prior without collapsing and preserving the base perplexity
of 3.8447.}
\label{fig:rlhf_curve}
\end{figure}

\subsubsection{Before vs. After RLHF Human Scores}

Table~\ref{tab:rlhf_results} summarises the aggregate human listening scores
before and after RLHF fine-tuning. The mean score improved from 3.7056
($\sigma = 0.9815$, $N=18$ participants, 180 total ratings) to 3.9941
($\sigma = 0.9671$, $N=17$ participants, 170 total ratings), yielding an
absolute improvement of $+0.2886$ points and a relative improvement of $7.79\%$
on the 1--5 Likert scale.

\begin{table}[htbp]
\caption{Aggregate Human Listening Scores Before and After RLHF (Task 4)}
\begin{center}
\begin{tabular}{|l|c|c|c|c|}
\hline
\textbf{Condition} & \textbf{Mean} & \textbf{Std} & \textbf{Participants} & \textbf{Ratings} \\
\hline
Before RLHF & 3.7056 & 0.9815 & 18 & 180 \\
After RLHF  & 3.9941 & 0.9671 & 17 & 170 \\
\hline
\textbf{Improvement} & \textbf{+0.2886} & -- & -- & -- \\
\hline
\end{tabular}
\label{tab:rlhf_results}
\end{center}
\end{table}

Fig.~\ref{fig:rlhf_comparison} shows the overall average scores before and
after RLHF (top panel) and the per-sample breakdown across all 10 evaluated
samples (bottom panel). The overall average rose from 3.7056 to 3.9941.
Nine out of ten samples improved after RLHF fine-tuning.

\begin{figure}[htbp]
\centerline{\includegraphics[width=\columnwidth]{architecture_diagram/task4_before_vs_after_comparison.png}}
\caption{Task 4 -- Before vs.\ After RLHF human listening scores. Top: overall
average score (Before RL = 3.7056, After RL = 3.9941). Bottom: per-sample mean
scores across 10 evaluated samples. Orange bars (After RL) exceed blue bars
(Before RL) in 9 out of 10 samples.}
\label{fig:rlhf_comparison}
\end{figure}

Table~\ref{tab:rlhf_persample} reports the per-sample scores before and after
RLHF fine-tuning. The largest improvement was observed in sample\_02
($+0.8889$), which was the lowest-scoring sample before RLHF (3.1111),
suggesting that the policy gradient updates particularly benefited weaker
generations. Samples with higher pre-RLHF scores generally showed smaller
gains, consistent with a ceiling effect. The only regression occurred in
sample\_10 ($-0.3333$), which held the highest pre-RLHF score (4.3333),
further supporting the presence of a score ceiling near 4.4 on this corpus.

\begin{table}[htbp]
\caption{Per-Sample Human Listening Scores Before and After RLHF (Task 4)}
\begin{center}
\begin{tabular}{|l|c|c|c|}
\hline
\textbf{Sample} & \textbf{Before RL} & \textbf{After RL} & \textbf{$\Delta$} \\
\hline
sample\_01 & 3.5000 & 3.7059 & +0.2059 \\
sample\_02 & 3.1111 & 4.0000 & +0.8889 \\
sample\_03 & 3.8333 & 3.8824 & +0.0490 \\
sample\_04 & 4.0000 & 4.4118 & +0.4118 \\
sample\_05 & 3.8333 & 4.1176 & +0.2843 \\
sample\_06 & 3.6667 & 4.0588 & +0.3922 \\
sample\_07 & 3.2778 & 3.8235 & +0.5458 \\
sample\_08 & 3.7778 & 3.9412 & +0.1634 \\
sample\_09 & 3.7222 & 4.0000 & +0.2778 \\
sample\_10 & 4.3333 & 4.0000 & $-$0.3333 \\
\hline
\end{tabular}
\label{tab:rlhf_persample}
\end{center}
\end{table}

As illustrated in Figs.~\ref{fig:before_rl} and~\ref{fig:after_rl} the individual
Before and after RLHF, sample-level average human scores are shown respectively. Before
RLHF, scores ranged from 3.1111 (sample\_02) to 4.3333 (sample\_10), with a
standard deviation of 0.9815. The post-RLHF score distribution became narrower and
shifted upward scores ranged from 3.7059 (sample\_01) to 4.4118 (sample\_04),
having a smaller standard deviation 0.9671, a larger floor and
better consistency in quality of samples.

\begin{figure}[htbp]
\centerline{\includegraphics[width=\columnwidth]{architecture_diagram/task4_before_rl_human_scores.png}}
\caption{Task 4 -- Per-sample average human listening scores before RLHF
fine-tuning (18 participants, 180 total ratings). Scores range from 3.1111
to 4.3333 with mean 3.7056.}
\label{fig:before_rl}
\end{figure}

\begin{figure}[htbp]
\centerline{\includegraphics[width=\columnwidth]{architecture_diagram/task4_after_rl_human_scores.png}}
\caption{Task 4 -- Per-sample average human listening scores after RLHF
fine-tuning (17 participants, 170 total ratings). Scores range from 3.7059
to 4.4118 with mean 3.9941, reflecting a 7.79\% relative improvement over
the pre-RLHF baseline.}
\label{fig:after_rl}
\end{figure}


% ============================================================
\section{Discussion}
\label{sec:discussion}

\subsection {Fidelity Diversity Trade-off and VAE Latent Interpolation}

Fig.~\ref{fig:all_metrics} uncovers an apparent fidelity--diversity trade-off. The
LSTM Autoencoder has the lowest Pitch Histogram Distance (0.1292) however.
Breaks to pure repetition of output (Repetition Ratio = 1.0000), a known.
failure mode of deterministic decoders. Its converged loss of MSE reconstruction.
of 0.0339 validates effective single-genre reconstruction on MAESTRO, but the
diversity is not possible in the absence of stochasticity. This is solved by the VAE through KL.
regularization, Rhythm Diversity = 0.4167 and Repetition Ratio =
0.1773 at the expense of higher pitch distance (0.5868); and optimum validation ELBO.
of 0.0217 (epoch 7) is the steady training period prior to KL collapse.
Latent interpolation between Classical and Jazz vectors generated genre-blended.
compositions, the VAE learns a smooth semantically structured.
latent space. The KL spike near epoch 18 ($\sim$$5.9\!\times\!10^9$) reflectss
uncontrolled $\beta$ growth; this would be inhibited by KL annealing in the future.

\subsection{Transformer Coherence and RLHF Quality Gains}

The Transformer has the best balanced profile: perplexity 3.8447,
Repetition Ratio 0.2298, and the Rhythm Diversity 0.1532, and train and validation.
curves intersecting by epoch 4, which shows a powerful generalisation on 2.79M token.
windows. It has higher Pitch Histogram Distance (0.7334) which is anticipated when modeling.
five genres simultaneously. The capacity of self-attention to attend all over the world.
long sequences is at the basis of its structural coherence~\cite{vaswani2017attention}.

The subjective listening quality cannot be optimized by maximum-likelihood training.
directly. This gap is closed by RLHF: the mean scores of human beings increased to 3.9941 
($+7.79\%$) on 20 steps of policy gradient, 9 out of 10 samples have improved.
Tasks 1--3 and the baselines were not evaluated in terms of human listening;
the rating of Transformer-generated is reflected in the before-RAHF rating score of 3.7056.
before fine-tuning. The Transformer is retained in the RLHF-tuned model.
base perplexity of 3.8447, since the updates made by the policy gradient are optimizing human.
preference and not token-level likelihood, and leaving the base perplexity.
unchanged. Constant sequence log-probability  ($\approx -475$) across
training verifies that policy did not fail. The sole regression
(sample\_10, $-0.3333$) is consistent with a score ceiling near 4.4 and
emphasizes the role of calibration of reward models in subsequent rounds.

\subsection{limitations and Future Work}

The main constraints are: LSTM AE the output of collapse, VAE KL instabilities, high.
Transformer Pitch distance, small RLHF participant population (N=1718),
and the failure of MIDI to reproduce the timbre or expressive dynamics. Future work
should use KL annealing, increase human evaluation sample, execute more.
The pairwise preference reward models~\cite{christiano2017deep} round with RLHF.
diffusion-based generation of symbolic music~\cite{mittal2021symbolic} explores
and instantiate a co-composition interface based on continuous feedback.

% ============================================================
\section{Deliverables}
\label{sec:deliverables}

The complete project submission includes the following artefacts, all organized
under the \texttt{outputs/} directory of the project GitHub repository.

\subsection*{A. Generated MIDI Files}
A total of 58 MIDI files are included across all tasks and baselines, far
exceeding the required minimum of 5--15 files. Table~\ref{tab:midi_files}
provides the complete breakdown.

\begin{table}[htbp]
\caption{Summary of All Generated MIDI Output Files}
\begin{center}
\begin{tabular}{|l|l|c|}
\hline
\textbf{Task} & \textbf{Description} & \textbf{\# Files} \\
\hline
Baseline & Random Note Generator        & 5  \\
Baseline & Markov Chain                 & 5  \\
Task 1   & LSTM AE --- Classical        & 5  \\
Task 2   & VAE --- Multi-Genre Samples  & 8  \\
Task 2   & VAE --- Latent Interpolation & 5  \\
Task 3   & Transformer Compositions     & 10 \\
Task 4   & RLHF Before Fine-Tuning      & 10 \\
Task 4   & RLHF After Fine-Tuning       & 10 \\
\hline
\textbf{Total} &                        & \textbf{58} \\
\hline
\end{tabular}
\label{tab:midi_files}
\end{center}
\end{table}

Task 1 samples represent Classical piano generation from the MAESTRO-trained
LSTM Autoencoder. Task 2 samples span Classical, Jazz, Rock and Pop genres just as the VAE model was trained on the Lakh MIDI dataset, which includes a variety of music genres. The model uses stochastic latent sampling from $z \sim \mathcal{N}(0, I)$; allowing it to generate music across these genres. Additionally, five interpolation files were generated, smoothly transitioning between Classical and Jazz latent vectors. Task 3 Transformer compositions are long-form autoregressive outputs
generated token-by-token across multiple genres. Task 4 files comprise the
complete before/after survey pairs used in the RLHF human listening evaluation.

\subsection*{B. Evaluation Results and Plots}
Quantitative evaluation outputs include metric summaries covering pitch
histogram distance ($\pm$std), rhythm diversity ($\pm$std), and repetition
ratio ($\pm$std) for all 38 model-generated MIDI files across Tasks 1--3 and
both baselines, along with machine-readable aggregate statistics per model
group. A Transformer perplexity report confirms the best validation perplexity
of 3.8447 and best validation loss of 1.3467, which is also retained by the
RLHF-tuned model. Full epoch-by-epoch training and validation loss histories
are included for Tasks 1--3, and RLHF training dynamics over 20 RL steps
including per-step loss and log-probability are also provided. Ten evaluation
plots are included covering training curves, perplexity, cross-task metric
comparisons, and before/after RLHF human score distributions.

\subsection*{C. Human Listening Survey Data}
The complete raw survey data is included for both evaluation rounds: 180 total
ratings from 18 participants before RLHF and 170 ratings from 17 participants
after RLHF. Long-format versions of both datasets are provided for statistical
analysis. Per-sample average reward scores used as RLHF reward signals are
included for both rounds, along with a per-sample comparison confirming that
9 of 10 samples improved after fine-tuning, aggregate statistics confirming a
$+0.2886$ absolute improvement ($+7.79\%$ relative), and summary reward
statistics for both survey rounds.

\subsection*{D. Trained Model Checkpoints}
All four trained model weights are saved and included. These are the best LSTM
Autoencoder checkpoint trained for 30 epochs on MAESTRO (Task 1, converged MSE
loss = 0.0339), the best VAE checkpoint saved from the stable training phase
prior to the KL explosion (Task 2, best validation ELBO = 0.0217 at epoch 7),
the best Transformer checkpoint achieving a best validation loss of 1.3467
(Task 3, 10 epochs), and the RLHF fine-tuned Transformer after 20 policy
gradient update steps retaining the base perplexity of 3.8447 (Task 4).


% ============================================================
\section{Conclusion}
\label{sec:conclusion}

In this work, we presented a four-stage progressive framework for unsupervised
multi-genre music generation. Beginning with an LSTM Autoencoder for
single-genre reconstruction, we introduced a VAE for diverse probabilistic
generation, a Transformer decoder for long-form sequence modeling, and an RLHF
stage to align generation with human preferences.

The LSTM Autoencoder achieved the lowest Pitch Histogram Distance (0.1292) but
collapsed to fully repetitive outputs (Repetition Ratio = 1.0000), with a
converged MSE reconstruction loss of 0.0339. The VAE improved Rhythm Diversity
to 0.4167 and reduced the Repetition Ratio to 0.1773, reaching a best
validation ELBO of 0.0217; latent interpolation between Classical and Jazz
latent vectors further demonstrated the model's capacity for smooth genre
transitions, with the resulting MIDI files included in the final submission.
The Transformer achieved a best validation perplexity of 3.8447 (loss = 1.3467)
over 2.79M token windows, with perplexity falling from 4.42 to 3.8447 on the
validation set. Human listening evaluation was conducted exclusively in Task 4,
with the pre-RLHF baseline score of 3.7056 reflecting Transformer-generated
sample ratings prior to fine-tuning. RLHF fine-tuning improved mean human
listening scores from 3.7056 to 3.9941 ($+7.79\%$, absolute improvement
$+0.2886$) across 18 participants, with 9 of 10 samples improving. The
RLHF-tuned model retains the Transformer base perplexity of 3.8447, as policy
gradient updates optimize for human preference rather than token-level
likelihood.

These results demonstrate that unsupervised deep generative models, augmented
with human preference alignment and latent-space interpolation, provide a
scalable approach to automatic music generation without manual annotation, with
clear potential for composition assistance, game audio, and adaptive music
systems.


% ============================================================
\section*{Acknowledgment}

The authors thank the instructors of CSE425/EEE474 Neural Networks for the
project guidance and course infrastructure. Computational experiments were
conducted using Google Colab with an NVIDIA GeForce RTX 4080 GPU.
The MAESTRO and Lakh MIDI datasets are gratefully acknowledged for enabling
this research.


% ============================================================
\bibliographystyle{IEEEtran}
\bibliography{reference.bib}

\end{document}